\documentclass[11pt,a4paper]{article}
\usepackage[T1]{fontenc}
\usepackage[utf8]{inputenc}
\usepackage[ngerman]{babel}
\usepackage[margin=2.5cm]{geometry}
\usepackage[hidelinks]{hyperref}
\usepackage{array}
\setlength{\parindent}{0pt}
\setlength{\emergencystretch}{2em}
\renewcommand{\arraystretch}{1.25}
\begin{document}
\shorthandoff{"}
\begin{center}
  {\Large\textbf{Webtechnologien -- Übungsblatt 3}}
\end{center}
\section*{Aufgabe 2: Making an HTML Page Editable}
\textbf{Aufgabe:} Zur Seite \texttt{editable.html} (drei
\texttt{section}s mit je einer \texttt{h1}) mit \texttt{editable.js} ein
Formular hinzufügen, mit dem man eine Section auswählen, ihre Überschrift
und ihr HTML bearbeiten und speichern kann. Am HTML darf nur die
Script-Referenz ergänzt werden.

\medskip
\textbf{Ergebnis:}
\begin{itemize}
  \item \texttt{editable.html}: Vorlage aus \texttt{editable.zip}, im
    \texttt{head} nur diese Zeile ergänzt:
    \begin{center}
    \texttt{<script src="editable.js" defer></script>}
    \end{center}
  \item Formular: \texttt{form}, \texttt{select}, Textfeld,
    \texttt{textarea} und Submit-Button werden mit
    \texttt{createElement} erzeugt und mit \texttt{after} direkt hinter die letzte Section gesetzt.
    Der Button ist zuerst \texttt{disabled}.
  \item Select: erste Option \glqq\texttt{-{}- choose a section -{}-}\grqq{}
    als Hinweis, danach pro Section eine Option. Der \texttt{value} ist
    die Nummer der Section, der Text die Überschrift aus
    \texttt{querySelector("h1")}.\allowbreak\texttt{textContent}.
  \item \texttt{change}-Listener am Select: Textfeld = Überschrift,
    Textarea = \texttt{section.innerHTML}, Button wird aktiv.
  \item \texttt{submit}-Listener am Formular:
    \texttt{event.preventDefault()} verhindert das Neuladen der Seite.
    Zuerst wird \texttt{section.innerHTML} auf den Inhalt der Textarea
    gesetzt. Dann wird die \texttt{h1} gesucht; fehlt sie, weil sie in der
    Textarea gelöscht wurde, wird eine neue erzeugt und mit
    \texttt{prepend} an den Anfang der Section gesetzt. Danach bekommt die
    \texttt{h1} den Text aus dem Textfeld. Das Textfeld gewinnt also immer
    gegen eine Änderung der Überschrift in der Textarea.
  \item Danach wird der Text der gewählten Option auf die neue Überschrift
    gesetzt und die Textarea zeigt den neuen Stand der Section.
\end{itemize}

\textbf{Hinweis:} Das Textfeld ist \texttt{required}, eine leere
Überschrift lässt der Browser nicht zu. Der Inhalt der Textarea wird
bewusst als HTML übernommen, weil die Angabe genau das verlangt; die
Änderung betrifft nur die eigene Seite im Browser, nicht die Datei.

\medskip
\textbf{Test:} In Chrome, \texttt{editable.html} direkt geöffnet:
\begin{center}
\begin{tabular}{|>{\raggedright\arraybackslash}p{5.6cm}|>{\raggedright\arraybackslash}p{8.1cm}|}
\hline
\textbf{Aktion} & \textbf{Ergebnis} \\ \hline
Seite laden & Formular nach der dritten Section, Optionen
  Topic 1, Topic 2, Topic 3, Button deaktiviert \\ \hline
\glqq Topic 2\grqq{} wählen & Textfeld \glqq Topic 2\grqq, Textarea mit
  \texttt{<h1>Topic 2</h1>...}, Button aktiv \\ \hline
In der Textarea \texttt{h1} auf \glqq FromTextarea\grqq{} und Text ändern,
  im Textfeld \glqq FromInput\grqq, Klick &
  Überschrift \glqq FromInput\grqq, geänderter Text übernommen,
  Option heißt \glqq FromInput\grqq \\ \hline
\glqq Topic 1\grqq{} wählen, \texttt{<h1>Topic 1</h1>} in der Textarea
  löschen, Textfeld \glqq New One\grqq, Klick &
  Section hat genau eine \texttt{h1} \glqq New One\grqq{} ganz vorne,
  Option aktualisiert \\ \hline
\end{tabular}
\end{center}
Keine Fehler in der Konsole, die Seite wird beim Speichern nicht neu
geladen.

\medskip
{\small \textbf{Hilfsmittel:} Claude Code (Claude Opus 5.5) für Erstellung, Prüfung und Dokumentation.}
\end{document}
